\newpage
\chapter{AMR: Real-Time Performance Diagnostics for Load Balancing}

\section{The Performance Diagnostics Landscape}

Performance optimization in high-performance computing has traditionally relied on post-mortem
analysis of execution traces collected during application runs.
This approach involves instrumenting applications to capture detailed performance data, which is
then analyzed offline by performance engineers and domain scientists.
The workflow typically consists of trace collection, data export, and subsequent analysis using
specialized tools and frameworks.

While this methodology has proven effective for smaller-scale systems, it faces fundamental
limitations when applied to large-scale HPC applications.
The volume of trace data grows dramatically with system size, making collection and analysis
increasingly expensive and unwieldy.
More critically, the temporal and spatial separation between data collection and analysis results
in the loss of crucial context that is essential for accurate performance diagnosis.

The conventional approach treats performance analysis as a separate phase that occurs after the
completion of the primary computation.
This separation creates an artificial barrier that prevents real-time feedback and adaptation,
limiting the ability to distinguish between fundamental algorithmic issues and transient
environmental factors that could be addressed through runtime adjustments.

\section{The Diagnostic Challenge}

The central challenge in performance optimization for large-scale HPC applications lies in
distinguishing between fundamental workload imbalance and environmental system noise.
BSP (Bulk Synchronous Parallel) applications are extremely sensitive to variability, where even
small perturbations can cascade into significant performance degradation.
However, the root causes of such variability are often opaque to conventional analysis methods.

Traditional trace analysis tools collect vast amounts of data but lack the contextual information
necessary to determine whether observed performance anomalies represent algorithmic deficiencies
or environmental factors.
A performance spike might indicate genuine load imbalance that requires algorithmic intervention,
or it might reflect transient system conditions such as network congestion, memory pressure, or
resource contention that could be addressed through runtime adaptation.

Without access to real-time hardware context and system state, performance engineers are forced
to make educated guesses about the root causes of performance issues.
This limitation is particularly acute for complex applications like Adaptive Mesh Refinement (AMR)
codes, where the dynamic nature of computational workloads makes it difficult to predict and
characterize performance behavior using static analysis alone.

\section{AMR Case Study and Investigation}

Our investigation focused on AMR applications as a representative example of the diagnostic
challenges facing modern HPC performance optimization.
AMR codes present particularly complex load balancing challenges due to their dynamic refinement
patterns, which create unpredictable computational workloads that evolve throughout the
application's execution.

The initial approach followed conventional methodology, collecting detailed execution traces and
attempting to identify load imbalance patterns through post-mortem analysis.
However, this approach quickly revealed the limitations of offline analysis when applied to
complex, dynamic workloads.
The traces captured performance symptoms but provided insufficient context to determine whether
observed imbalances were fundamental to the algorithmic structure or could be addressed through
runtime adjustments.

This limitation drove the development of a tunable placement policy designed to experimentally
manipulate load distribution and observe the resulting performance characteristics.
The policy enabled controlled experiments that could isolate the effects of different load
balancing strategies, providing insights that would be impossible to obtain through static
analysis alone.

\section{In-Situ Diagnostic Approach}

The inadequacy of conventional analysis methods necessitated the development of new approaches
that could provide real-time diagnostic capabilities with full hardware context.
The investigation evolved toward on-site, full-stack performance analysis that could capture
both application-level behavior and system-level conditions simultaneously.

This approach required the development of lightweight monitoring infrastructure that could collect
targeted performance data without significantly impacting the primary computation.
The system needed to provide programmable data collection capabilities that could adapt to
different diagnostic scenarios while maintaining minimal overhead.

The work organically converged toward a telemetry pipeline architecture using OLAP databases and
SQL for performance data analysis.
This approach proved more effective than traditional trace formats, enabling complex analytical
queries that could correlate application behavior with system conditions.
However, the limitations of post-hoc analysis remained, highlighting the need for real-time
feedback capabilities.

\subsection{Programmable Telemetry Pipeline}

The development of programmable telemetry capabilities represented a significant departure from
conventional performance analysis workflows.
Rather than relying on predefined instrumentation points and fixed data collection strategies,
the system enabled dynamic adaptation of monitoring behavior based on observed conditions.

The pipeline architecture supported complex analytical queries that could identify performance
patterns and correlate them with system state in real-time.
This capability proved essential for distinguishing between fundamental load imbalance and
environmental factors that could be addressed through runtime intervention.

\subsection{eBPF Integration and Insights}

The investigation's exploration of eBPF technologies provided crucial insights into the potential
for programmable, real-time performance monitoring.
eBPF's ability to deploy custom monitoring logic directly within the kernel enabled unprecedented
visibility into system behavior without the overhead and complexity of traditional instrumentation
approaches.

While eBPF provided powerful per-process collection capabilities, the challenge of aggregating
and analyzing this data in real-time remained.
The experience highlighted the need for integrated frameworks that could seamlessly combine
programmable data collection with sophisticated analytical capabilities.

\section{Experimental Validation and Results}

The experimental validation of the in-situ diagnostic approach demonstrated significant advantages
over conventional post-mortem analysis methods.
The real-time feedback capabilities enabled rapid identification of performance issues and
immediate validation of proposed solutions.

The tunable placement policy experiments revealed that many apparent load imbalance problems were
actually manifestations of environmental factors that could be addressed through runtime
adaptation.
This finding fundamentally challenged the conventional assumption that performance optimization
requires algorithmic changes and established the importance of real-time system context for
accurate diagnosis.

Performance improvements achieved through the in-situ approach significantly exceeded those
obtained through conventional optimization methods.
The ability to distinguish between fundamental and environmental performance issues enabled
targeted interventions that addressed root causes rather than symptoms.

\section{Lessons and Implications}

The AMR investigation established that performance optimization requires the same real-time
observation and control capabilities that proved essential for data management in the CARP work.
The fundamental pattern of continuous monitoring, analysis, and adaptation applies equally to
both domains, suggesting that these apparently distinct challenges are manifestations of the
same underlying requirement.

The work demonstrated that performance analysis must happen in the hardware context where the
application is executing.
The loss of temporal and spatial locality inherent in conventional post-mortem analysis
eliminates crucial information that is essential for accurate diagnosis.
This finding has profound implications for the design of performance optimization tools and
methodologies.

Perhaps most importantly, the AMR work revealed that effective performance optimization requires
programmable analytical capabilities that can adapt to changing conditions and requirements.
The emergence of SQL-based analytical workflows and the insights gained from eBPF integration
pointed toward the need for unified frameworks that could provide both data collection and
analysis capabilities within a single, coherent system.

The convergence of lessons from both the data management domain (CARP) and the performance
optimization domain (AMR) establishes the foundation for a unified approach to in-situ
observation and control that transcends traditional domain boundaries. 